\documentclass[11pt,a4paper]{article}

% --- Encodage & Langue ---
\usepackage[utf8]{inputenc}
\usepackage[T1]{fontenc}

\usepackage[french]{babel}

% Renommage manuel en français (si babel n'est pas activé)
\renewcommand{\contentsname}{Table des matières}

% --- Mise en page & Typographie ---
\usepackage[margin=2.5cm]{geometry}
\usepackage{parskip} % Espacement entre paragraphes sans alinéa
\setlength{\headheight}{15pt}

% Polices optionnelles (décommenter si texlive-fontsrecommended est installé) :
\usepackage{lmodern}
\usepackage{microtype}

% --- Couleurs & Liens hypertextes ---
\usepackage{xcolor}
\usepackage{hyperref}
\definecolor{linkblue}{RGB}{0, 50, 180}
\hypersetup{
    colorlinks=true,
    linkcolor=linkblue,
    urlcolor=linkblue
}

% --- En-têtes et pieds de page ---
\usepackage{fancyhdr}
\pagestyle{fancy}
\fancyhf{}
\lhead{LAN Corpauration -- Consignes \& Règlement}
\rhead{\leftmark}
\cfoot{\thepage}

% --- Informations du document ---
\title{\textbf{LAN -- Consignes \& Informations Générales}}
\author{Corpauration}
% \date{\today}

\begin{document}

    \maketitle
    \thispagestyle{fancy}

    \tableofcontents
    \vspace{1.5em}
    %\hrule
    %\vspace{1.5em}
    \newpage


    \section{Informations Pratiques}
    \subsection{Contacts}
    \begin{itemize}
        \item Discord : \url{https://discord.gg/Z9cnuuChjc}
        \item Mail : \url{asso.corpauration@gmail.com}
    \end{itemize}


    \section{Matériel}

    \subsection{Équipement Obligatoire}
    \begin{itemize}
        \item Ordinateur et 1 seul écran
        \item Périphériques (clavier, souris, câbles, etc.)
        \item Casque audio et/ou écouteurs (\textit{enceintes interdites})
        \item Câble Ethernet RJ45
        \item Multiprise
    \end{itemize}

    \subsection{Préparation en amont}
    \begin{itemize}
        \item Effectuer les mises à jour système
        \item Télécharger et mettre à jour \textbf{TOUS} les jeux du tournoi à l'avance
    \end{itemize}

    \section{Réseau \& Connexion}

    \subsection{Configuration réseau}
    \begin{itemize}
        \item Se connecter en filaire à une prise réseau.
        \item Se rendre à l'adresse \url{https://lan.corpauration.fr}.
        \item Aller voir les DSI en salle 213 pour confirmer la connexion.

    \end{itemize}

    \subsection{Pas de téléchargements massifs durant les tournois}
    Afin de préserver la bande passante et la stabilité du réseau local pour tous les participants, tout téléchargement lourd, streaming personnel ou mise à jour non planifiée pendant les épreuves est strictement interdit.
    Bien penser à télécharger tous les jeux avant la LAN\@.


    \section{Plateforme Web du Tournoi}

    Une plateforme web locale dédiée au tournoi sera disponible pendant l'évènement.

    \subsection{Classement en Direct}
    Les participants peuvent suivre en temps réel l'évolution des scores, des classements par épreuve et du classement général grâce à une synchronisation instantanée.

    \subsection{Espace Capitaine \& Remontée des Scores}
    Chaque capitaine d'équipe dispose d'un accès dédié pour gérer les rencontres de son équipe :
    \begin{itemize}
        \item \textbf{Connexion :} Authentification avec les identifiants remis lors de l'accueil.
        \item \textbf{Canal exclusif :} La plateforme constitue le \textbf{seul et unique moyen officiel} de transmettre les résultats.
        Aucune déclaration à l'oral, sur Discord ou par messagerie ne sera prise en compte.
        \item \textbf{Preuve obligatoire :} Le capitaine doit obligatoirement joindre une capture d'écran nette du tableau de fin de partie (import de fichier ou collage direct via le raccourci \texttt{Ctrl+V}).
        \item \textbf{Résultats non conformes :} Tout résultat absent, transmis hors délai ou incomplet (absence de capture, image tronquée ou illisible) sera automatiquement sanctionné par un \textbf{score de 0 point}.
        \item \textbf{Arbitrage :} Les organisateurs et arbitres vérifient les preuves avant validation définitive au classement général.
    \end{itemize}


    \section{Déroulement des Tournois}

    \subsection{Système de Points Général}
    À l'issue de chaque jeu, un classement d'épreuve est établi.
    Les équipes reçoivent des points pour le classement général selon un barème dégressif annoncé en début de tournoi.\\
    Par exemple, pour 10 équipes le tableau est le suivant :

    \begin{center}
        \begin{tabular}{|c|c||c|c|}
            \hline
            \textbf{Rang} & \textbf{Points} & \textbf{Rang} & \textbf{Points} \\
            \hline
            1\ier{}       & 25 pts          & 6\ieme{}      & 8 pts           \\
            2\ieme{}      & 18 pts          & 7\ieme{}      & 6 pts           \\
            3\ieme{}      & 15 pts          & 8\ieme{}      & 4 pts           \\
            4\ieme{}      & 12 pts          & 9\ieme{}      & 2 pts           \\
            5\ieme{}      & 10 pts          & 10\ieme{}     & 1 pt            \\
            \hline
        \end{tabular}
    \end{center}
    L'équipe cumulant le plus de points sur l'ensemble des 6 épreuves remporte le tournoi.

    \subsection{Règlement spécifique par jeu}

    \subsubsection{Apex Legends}
    \textbf{Mode :} Salon personnalisé en Trios.

    \paragraph{Format de l'épreuve}
    \begin{itemize}
        \item \textbf{Durée allouée prévisionnelle :} 4 heures.
        Les parties s'enchaînent dans le créneau imparti.
        \item \textbf{Accès à la partie :} Le code pour rejoindre le lobby est envoyé par les administrateurs à tous les capitaines.
        \item \textbf{Composition de l'escouade :} Chaque trio doit obligatoirement intégrer au moins une légende parmi celles disponibles par défaut pour un nouveau joueur (Bloodhound, Gibraltar, Lifeline, Pathfinder, Wraith, Bangalore).
    \end{itemize}

    \paragraph{Calcul du score}
    Le score de chaque manche suit le barème officiel compétitif (ALGS) :
    \begin{itemize}
        \item \textbf{Éliminations :} +1 point par frag d'équipe.
        \item \textbf{Points de placement par manche :}
    \end{itemize}

    \begin{center}
        \begin{tabular}{|c|c||c|c|}
            \hline
            \textbf{Placement} & \textbf{Points} & \textbf{Placement} & \textbf{Points} \\
            \hline
            1\ier{}            & 12 pts          & 6\ieme{} à 7\ieme{}   & 3 pts           \\
            2\ieme{}           & 9 pts           & 8\ieme{} à 10\ieme{}  & 2 pts           \\
            3\ieme{}           & 7 pts           & 11\ieme{} à 15\ieme{} & 1 pt            \\
            4\ieme{}           & 5 pts           & 16\ieme{} à 20\ieme{} & 0 pt            \\
            5\ieme{}           & 4 pts           & Frag                   & 1 pt / élim.    \\
            \hline
        \end{tabular}
    \end{center}

    \textit{Le cumul des points (placement + frags) sur toutes les parties de la session détermine le classement de l'épreuve.}

    \subsubsection{League of Legends [À vérifier]}
    \textbf{Mode :} ARAM du chaos.

    \paragraph{Déroulement}
    \begin{itemize}
        \item \textbf{Structure :} Affrontements 2 à 2 par rondes avec cumul de points sous forme de championnat suivi d'une finale.
        \item \textbf{Format des matchs :} Match en 1 manche gagnante.
        \item \textbf{Attribution des champions :} Sélection parmi les champions choisis aléatoirement par le jeu.
        \item \textbf{Barème par match :}
        \begin{itemize}
            \item \textbf{Victoire :} +3 points
            \item \textbf{Défaite :} +1 point
            \item \textbf{Forfait / Non-présentation :} 0 point
        \end{itemize}
    \end{itemize}

    \paragraph{Condition de victoire} Destruction du Nexus adverse.

    \subsubsection{Fall Guys}
    \textbf{Mode :} Tournoi Solo en salon personalisé.

    \paragraph{Calcul du score d'équipe}
    Le score de l'équipe est calculé par la somme des performances de ses 3 membres :
    \begin{itemize}
        \item \textbf{Survie :} +1 point par manche franchie par joueur.
        \item \textbf{Prime de victoire :} Bonus accordé en cas de victoire de la partie de 3pts.
    \end{itemize}

    \paragraph{Classement final}
    Somme des points totaux obtenus par l'équipe sur toute la durée de l'épreuve.

    \subsubsection{GeoGuessr}
    \textbf{Mode :} Challenge d'équipe.

    \paragraph{Format de l'épreuve}
    \begin{itemize}
        \item Une série sur chacune des cartes sélectionnées.
        \item Mode Solo, No Move, 1 min 15 par round.
        \item Coopération des 3 membres possible, mais sans aide extérieure (\textit{Moteurs de recherche, IA, guides, etc.}, interdits).
    \end{itemize}

    \paragraph{Calcul du score}
    À l'issue de chaque manche (carte), les équipes sont classées selon leur score brut GeoGuessr et reçoivent des points de placement selon le barème général :

    \begin{center}
        \begin{tabular}{|c|c||c|c|}
            \hline
            \textbf{Placement} & \textbf{Points} & \textbf{Placement} & \textbf{Points} \\
            \hline
            1\ier{}  & 25 pts & 6\ieme{}  & 8 pts \\
            2\ieme{} & 18 pts & 7\ieme{}  & 6 pts \\
            3\ieme{} & 15 pts & 8\ieme{}  & 4 pts \\
            4\ieme{} & 12 pts & 9\ieme{}  & 2 pts \\
            5\ieme{} & 10 pts & 10\ieme{} & 1 pt  \\
            \hline
        \end{tabular}
    \end{center}

    \textit{Le cumul des points de placement sur l'ensemble des manches détermine le classement final du jeu. En cas d'égalité, le score GeoGuessr brut cumulé sert de départage.}

    \paragraph{Cartes possibles}
    \begin{itemize}
        \item \href{https://www.geoguessr.com/fr/maps/world}{World}
        \item \href{https://www.geoguessr.com/fr/maps/62a44b22040f04bd36e8a914}{A Community World}
        \item \href{https://www.geoguessr.com/fr/maps/famous-places}{Famous Places}
        \item \href{https://www.geoguessr.com/fr/maps/france}{France}
        \item \href{https://www.geoguessr.com/fr/maps/64edb149ef5a06cfe58bf922}{F1 tracks 2026}
        \item \href{https://www.geoguessr.com/fr/maps/6554fb70c974b236f72cd734}{Pau (Intérieur de la rocade)}
        \item \href{https://www.geoguessr.com/fr/maps/60451069abd02500012138c9}{Eiffel Tower replicas and derivatives}
    \end{itemize}

    \subsubsection{Super Smash Bros Ultimate}
    \textbf{Mode :} Smash en bande en 3 contre 3.\\
    \textbf{Règles :}
    \begin{itemize}
        \item 3 vies
        \item Pas de smash final, d'objets, de dégâts initiaux
        \item Carte aléatoire en mode CB parmi les suivantes :
        \begin{itemize}
            \item Battlefield
            \item Small Battlefield
            \item Final Destination
            \item Yoshi's Story
            \item Pokemon Stadium 2
            \item SmashVille
            \item Lylat Cruise
            \item Hollow Bastion
            \item Town and City
        \end{itemize}
    \end{itemize}

    \paragraph{Format du tournoi}
    Le tournoi se passe sous forme de championnat.
    Toutes les équipes affronteront un certain nombre d'équipes et gagneront des points en fonction de leur victoire ou défaite.
    Les deux équipes avec le plus de points s'affrontent pour déterminer le gagnant du jeu (mêmes consignes, mais sous forme de BO3).

    \paragraph{Déroulement des duels}
    Chaque duel se déroule en une manche (hors finale entre les 2 meilleures équipes en BO3).
    Pour chaque round de ce duel :
    \begin{itemize}
        \item \textbf{Bannissements :} Chaque équipe bannit 3 personnages à tour de rôle (la première équipe à bannir est l'équipe perdante d'un \textit{Pierre Feuille Ciseaux} en une manche).
        \item \textbf{Sélection de la carte :} La map est sélectionnée aléatoirement par le jeu parmi la liste autorisée.
        \item \textbf{Ordre des affrontements :} L'ordre de passage des joueurs de chaque équipe est déterminé aléatoirement par la fonction dédiée dans le jeu (touche \texttt{X}).
    \end{itemize}

    \subsubsection{Mario Kart 8 Deluxe}
    \textbf{Mode :} Lobbies en mode LAN (12 joueurs), courses en format \textit{pseudo-solo}.

    \paragraph{Déroulement}
    \begin{itemize}
        \item \textbf{Configuration :} Salons créés via le mode \textit{LAN} regroupant jusqu'à 12 joueurs en simultané.
        \item \textbf{Format des rencontres :} Chaque équipe dispute 2 à 3 sessions (lobbies) contre des adversaires différents.
        \item \textbf{Format d'une session :} Session de 4 courses consécutives.
        \item \textbf{Calcul du score :} Chaque joueur concourt en individuel.
        \item \textbf{Résultats :} Les points récoltés individuellement par les 3 membres d'une équipe selon le barème du jeu sont cumulés sur l'ensemble de leurs parties pour établir le classement final de l'épreuve.
    \end{itemize}

    \paragraph{Paramètres de partie}
    \begin{itemize}
        \item Mode : Course VS (150cc)
        \item Règles d'équipe : Individuel / Chacun pour soi (\textit{pseudo-solo})
        \item Objets : Normaux
        \item IA : Sans IA
        \item Véhicules : Tous les véhicules autorisés
        \item Nombre de courses : 4 courses par session
    \end{itemize}

\end{document}
